\textbf{Proof.}
Fix \(\mathbb P_{n,m}\in\mathcal M_{n,m}(a,s,t)\), and put
\[
\boldsymbol\phi(u)=(1,u,\ldots,u^p)^\top,
\quad
\overline G_h=\int_0^1\boldsymbol\phi(u)\boldsymbol\phi(u)^\top u^a g_S(hu)\,du,
\quad
\overline q_h=\int_0^1\boldsymbol\phi(u)g_T(hu)\,du.
\tag{1}
\]
The change of variables \(x=hu\), the identities \(f_S(x)=x^ag_S(x)\) and \(f_T(x)=g_T(x)\), and the definition of \(\mathbf R_h\) give
\[
G_h^o=h^{a+1}\overline G_h,
\qquad
q_h^o=h\overline q_h.
\tag{2}
\]
For every nonzero \(c\in\mathbb R^{p+1}\),
\[
c^\top\overline G_hc
=\int_0^1\{c^\top\boldsymbol\phi(u)\}^2u^ag_S(hu)\,du>0.
\tag{3}
\]
Indeed, \(g_S(hu)\ge c_{S,-}>0\), \(u^a>0\) on \((0,1]\), and a nonzero polynomial cannot vanish on a set of positive Lebesgue measure. Thus \(\overline G_h\), and hence \(G_h^o\), is positive definite. This proves every inverse used below exists.

Let
\[
H_a=\int_0^1\boldsymbol\phi(u)\boldsymbol\phi(u)^\top u^a\,du.
\]
The same polynomial argument shows that \(H_a\) is positive definite. The source-factor bounds imply the Loewner inequalities
\[
c_{S,-}H_a\preceq\overline G_h\preceq c_{S,+}H_a.
\tag{4}
\]
Moreover, the first coordinate of \(\overline q_h\) lies in \([c_{T,-},c_{T,+}]\), and all its coordinates are bounded in absolute value by \(c_{T,+}\). Since \(p\) is fixed, (4) therefore yields constants \(0<c<C<\infty\), independent of the law and of sufficiently small \(h\), such that
\[
c\le \overline q_h^\top\overline G_h^{-1}\overline q_h\le C.
\tag{5}
\]
Combining (2) and (5) proves
\[
(q_h^o)^\top(G_h^o)^{-1}q_h^o
=h^{1-a}\overline q_h^\top\overline G_h^{-1}\overline q_h
\asymp h^{1-a}.
\tag{6}
\]

If \(\pi(x)=d^\top\mathbf R_h(x)\) has degree at most \(p\), then the definition of \(w_h^{\mathrm{BLP}}\) gives
\[
\begin{aligned}
\int_0^h w_h^{\mathrm{BLP}}(x)\pi(x)f_S(x)\,dx
&=d^\top G_h^o(G_h^o)^{-1}q_h^o\\
&=d^\top q_h^o
=\int_0^h\pi(x)f_T(x)\,dx.
\end{aligned}
\tag{7}
\]
This proves exact polynomial reproduction, including the constant coordinate.

Equations (1)--(2) also give, for \(0\le u\le1\),
\[
w_h^{\mathrm{BLP}}(hu)
=h^{-a}\boldsymbol\phi(u)^\top\overline G_h^{-1}\overline q_h.
\tag{8}
\]
The eigenvalue bounds in (4), boundedness of \(\overline q_h\), and boundedness of \(\boldsymbol\phi\) on \([0,1]\) imply
\[
\int_0^h|w_h^{\mathrm{BLP}}(x)|f_S(x)\,dx
=h\int_0^1
\left|\boldsymbol\phi(u)^\top\overline G_h^{-1}\overline q_h\right|
u^ag_S(hu)\,du
\le Ch.
\tag{9}
\]
Also, by direct multiplication,
\[
\int_0^h\{w_h^{\mathrm{BLP}}(x)\}^2f_S(x)\,dx
=(q_h^o)^\top(G_h^o)^{-1}q_h^o.
\tag{10}
\]
The two summands defining \(\omega_h^o\) have disjoint supports up to the Lebesgue-null point \(h\), while \(r^2f_S=f_T^2/f_S\) on \((0,1]\). Consequently (10) proves the exact identity
\[
\|\omega_h^o\|_{L^2(P_S)}^2
=(q_h^o)^\top(G_h^o)^{-1}q_h^o
+\int_h^1\frac{f_T(x)^2}{f_S(x)}\,dx.
\tag{11}
\]

We next bound the boundary bias. Write \(p=\lceil s\rceil-1\) and \(\beta=s-p\in(0,1]\). Under the stated Hölder convention, Taylor's formula at zero gives, for every \(v\in\mathcal H^s(2M)\),
\[
T_pv(x)=\sum_{j=0}^p\frac{v^{(j)}(0)}{j!}x^j,
\qquad
|v(x)-T_pv(x)|\le C_sM x^s,
\quad 0\le x\le1.
\tag{12}
\]
For \(p=0\), (12) is exactly the \(s\)-Hölder inequality. For \(p\ge1\), it follows from
\[
v(x)-T_pv(x)
=\frac1{(p-1)!}\int_0^x(x-y)^{p-1}
\{v^{(p)}(y)-v^{(p)}(0)\}\,dy
\]
and the \(\beta\)-Hölder bound on \(v^{(p)}\). By (7), the signed measure in \(b_h(v;P)\) annihilates \(T_pv\). Equations (9), (12), and \(\int_0^hf_T\le c_{T,+}h\) therefore imply
\[
\begin{aligned}
|b_h(v;P)|
&=\left|\int_0^h\{w_h^{\mathrm{BLP}}f_S-f_T\}(v-T_pv)\,dx\right|\\
&\le C h^s
\left\{\int_0^h|w_h^{\mathrm{BLP}}|f_S\,dx+
\int_0^hf_T\,dx\right\}
\le Ch^{s+1}.
\end{aligned}
\tag{13}
\]
Taking the supremum over the defining class of \(v\) proves \(b_h^o(P)\le Ch^{s+1}\). The response contrast \(\mu=\mu_1-\mu_0\) belongs to \(\mathcal H^s(2M)\) by the two response-smoothness assumptions and satisfies \(\|\mu\|_\infty\le1\) because \(0\le\mu_d\le1\). Hence
\[
|b_h(\mu;P)|\le b_h^o(P).
\tag{14}
\]
Notice also the exact oracle identity
\[
\int_0^1\omega_h^o(x)v(x)f_S(x)\,dx-
\int_0^1v(x)f_T(x)\,dx=b_h(v;P),
\tag{15}
\]
which distinguishes the actual bias in (14) from its worst-case allowance.

We now prove the sharpness assertion for \(a>0\). Use the admissible covariate densities
\[
f_S(x)=(a+1)x^a,
\qquad f_T(x)=1.
\tag{16}
\]
Their factors are the constants \(a+1\) and \(1\), so the boundary bounds, density smoothness conditions, and density normalizations hold when \(M\ge a+1\). Complete the full-data law by drawing \(A\) according to the known propensity, setting \(\mu_0(x)=1/2\) and \(\mu_1(x)=1/2+\delta x^s\), taking conditionally independent Bernoulli potential outcomes with these means in both populations, and imposing consistency. Choose \(\delta>0\) small enough that \(\delta\le\sqrt{1/4-M^{-1}}\) and that the two response functions lie in their radius-\(M\) Hölder balls. Treatment randomization and mean transportability then hold, and
\[
\operatorname{Var}(Y\mid A=d,X=x)\ge\frac14-\delta^2\ge M^{-1}.
\]
Thus the resulting observed law lies in \(\mathcal M_{n,m}(a,s,t)\), and its response contrast is the witness \(v_*(x)=\delta x^s\) used below.

Index rows and columns from zero and define
\[
A_{ij}=\frac1{a+i+j+1},
\qquad c_i=\frac1{i+1},
\qquad \boldsymbol d=A^{-1}\boldsymbol c,
\qquad P_a(u)=\sum_{j=0}^pd_ju^j.
\tag{17}
\]
The positive-definiteness argument in (3) shows that \(A\) is invertible. Under (16), scaling the defining equations for \(w_h^{\mathrm{BLP}}\) gives
\[
w_h^{\mathrm{BLP}}(hu)f_S(hu)h\,du
=hP_a(u)u^a\,du.
\tag{18}
\]
The preceding choice of \(\delta\), reduced further if necessary, ensures that \(v_*(x)=\delta x^s\) belongs to \(\mathcal H^s(2M)\) and has supremum norm at most one. Such a choice is possible because the \(p\)-th derivative of \(x^s\) is a constant multiple of \(x^\beta\), whose \(\beta\)-Hölder seminorm on \([0,1]\) is finite. From (18),
\[
b_h(v_*;P)=\delta h^{s+1}F_a(s),
\qquad
F_a(z)=\sum_{j=0}^p\frac{d_j}{a+z+j+1}-\frac1{z+1}.
\tag{19}
\]
The equations \(A\boldsymbol d=\boldsymbol c\) say that \(F_a(0)=\cdots=F_a(p)=0\). With
\[
D_a(z)=(z+1)\prod_{j=0}^p(z+a+j+1),
\]
the product \(D_aF_a\) is a polynomial of degree at most \(p+1\) with these \(p+1\) zeros. Evaluating at \(z=-1\), which is not a pole of the first term when \(a>0\), determines its constant and yields
\[
F_a(s)=
\frac{(-1)^p
\displaystyle\prod_{j=0}^p(a+j)
\displaystyle\prod_{j=0}^p(s-j)}
{(p+1)!(s+1)
\displaystyle\prod_{j=0}^p(s+a+j+1)}.
\tag{20}
\]
Because \(a>0\) and \(s>p\), (20) is nonzero. Equations (19)--(20) therefore give
\[
b_h^o(P)\ge |b_h(v_*;P)|\ge c h^{s+1},
\tag{21}
\]
which proves the class-sharp lower bound.

At the endpoint \(a=0\), the leading constant in (20) vanishes, and the claimed improvement follows from density regularity. Put \(\alpha=\min\{t,1\}\). The two density-factor assumptions imply, uniformly over the model class,
\[
|g_S(hu)-g_S(0)|+|g_T(hu)-g_T(0)|\le Ch^\alpha,
\qquad 0\le u\le1.
\tag{22}
\]
Let \(H_0=\int_0^1\boldsymbol\phi\boldsymbol\phi^\top\) and \(q_0=\int_0^1\boldsymbol\phi\). Since \(H_0e_0=q_0\), (22), the factor lower bound, and finite-dimensional inverse perturbation give
\[
\overline G_h^{-1}\overline q_h
=\frac{g_T(0)}{g_S(0)}e_0+O(h^\alpha)
\tag{23}
\]
uniformly. Equations (8), (22), and (23) then show
\[
\sup_{0\le x\le h}
|w_h^{\mathrm{BLP}}(x)f_S(x)-f_T(x)|\le Ch^\alpha.
\tag{24}
\]
Combining (12), polynomial reproduction, and (24) proves
\[
\sup_{\mathbb P_{n,m}\in\mathcal M_{n,m}(0,s,t)}b_h^o(P)
\le Ch^{s+1+\alpha}.
\tag{25}
\]
In the special case \(a=0\) and \(f_S=f_T\), one has \(G_h^oe_0=q_h^o\), whence \(w_h^{\mathrm{BLP}}=1\) on \([0,h)\) and \(b_h^o(P)=0\). Thus neither the actual bias nor its per-law envelope is being asserted equivalent to \(h^{s+1}\).

It remains to evaluate the deterministic allowance. Equations (6), (11), and \(\mathrm{lem:riesz\text{-}threshold}\) imply, uniformly over the model class,
\[
\|\omega_h^o\|_{L^2(P_S)}^2\asymp
\begin{cases}
1,&a<1,\\
\log(1/h),&a=1,\\
h^{1-a},&a>1.
\end{cases}
\tag{26}
\]
For \(a<1\), \(h_n=n^{-1/\{2(s+1)\}}\), so \(h_n^{s+1}=n^{-1/2}\), and (26) makes the oracle source-standard-deviation term of order \(n^{-1/2}\). For \(a=1\),
\[
h_n^{s+1}=\sqrt{\frac{\log n}{n}},
\qquad
\log(1/h_n)=\frac{\log n-\log\log n}{2(s+1)}\asymp\log n,
\tag{27}
\]
so both allowance terms have order \(\sqrt{\log n/n}\). For \(a>1\), with \(D=2s+a+1\),
\[
h_n^{s+1}=n^{-(s+1)/D},
\qquad
n^{-1/2}h_n^{(1-a)/2}
=n^{-1/2+(a-1)/(2D)}
=n^{-(s+1)/D}.
\tag{28}
\]
Equations (26)--(28) and \(\mathrm{def:frontier\text{-}rate}\) prove
\[
m^{-1/2}+h_n^{s+1}
+n^{-1/2}\|\omega_{h_n}^o\|_{L^2(P_S)}
\asymp m^{-1/2}+\rho_n(a,s).
\tag{29}
\]
The construction uses the unknown population densities through \(G_h^o,q_h^o\), and \(r\). Hence (29) is a deterministic oracle benchmark and supplies neither a data-measurable estimator nor a minimax upper bound. \(\square\)